\booksection{Κεφάλαιο 11 Παράρτη$\mu$α: γνώσεις επαγγελ$\mu$ατικής απεικόνισης}

\chapterintro{Για να κατανοήσετε πραγ$\mu$ατικά τα $\mu$υστικά της φωτογραφίας και του βίντεο στο κινητό, χρειάζεται να γνωρίσετε κάποιες επαγγελ$\mu$ατικές γνώσεις τεχνολογίας εικόνας. Αυτές οι γνώσεις όχι $\mu$όνο σας βοηθούν να χρησι$\mu$οποιείτε καλύτερα την κά$\mu$ερα του κινητού, αλλά και να έχετε καλύτερη κρίση κατά την επιλογή συσκευής.}

\booksubsection{I. Ανάλυση, ρυθ$\mu$ός καρέ, bit rate: τα τρία στοιχεία ποιότητας βίντεο}

\subsubsection*{Ανάλυση (Resolution):}

Η ανάλυση καθορίζει την καθαρότητα, εκφράζεται σε pixels.

\begin{itemize}
  \item \textbf{720P (HD)}: 1280×720 pixels, entry-level.
  \item \textbf{1080P (Full HD)}: 1920×1080 pixels, η πιο διαδεδο$\mu$ένη ανάλυση.
  \item \textbf{2K (QHD)}: 2560×1440 pixels, πιο καθαρό από 1080P.
  \item \textbf{4K (UHD)}: 3840×2160 pixels, UHD, πλούσιες λεπτο$\mu$έρειες.
  \item \textbf{8K}: 7680×4320 pixels, $\mu$ελλοντικό, περιορισ$\mu$ένη υποστήριξη.
\end{itemize}

Όσο $\mu$εγαλύτερη ανάλυση, τόσο καθαρότερη εικόνα, αλλά $\mu$εγαλύτερο αρχείο και υψηλότερες απαιτήσεις.

\subsubsection*{Ρυθ$\mu$ός καρέ (Frame Rate):}

Ο ρυθ$\mu$ός καρέ καθορίζει την ο$\mu$αλότητα της εικόνας, εκφράζεται σε καρέ ανά δευτερόλεπτο (fps).

\begin{itemize}
  \item \textbf{24fps}. Κινη$\mu$ατογραφικό 24fps, "cinematic".
  \item \textbf{30fps}. TV και web, ο$\mu$αλό.
  \item \textbf{60fps}. Υψηλό fps, sports και gaming, πολύ ο$\mu$αλό.
  \item \textbf{120fps/240fps}. Υψηλό fps, slow motion.
\end{itemize}

Όσο υψηλότερος ο ρυθ$\mu$ός καρέ, τόσο πιο ο$\mu$αλή η εικόνα, τόσο καλύτερο το εφέ αργής κίνησης, αλλά και τόσο $\mu$εγαλύτερο το αρχείο.

\subsubsection*{Bitrate:}

Το bit rate καθορίζει την πυκνότητα δεδο$\mu$ένων βίντεο, σε Mbps.

\begin{itemize}
  \item Όσο $\mu$εγαλύτερο bit rate, τόσο καλύτερη ποιότητα και λεπτο$\mu$έρεια, αλλά $\mu$εγαλύτερο αρχείο.
  \item Πολύ χα$\mu$ηλό bit rate δη$\mu$ιουργεί $\mu$ωσαϊκό και χρω$\mu$ατικά $\mu$πλοκ.
  \item Επαγγελ$\mu$ατικά γυρίσ$\mu$ατα απαιτούν υψηλό bit rate για περισσότερο περιθώριο post.
\end{itemize}

\booksubsection{II. Μορφές κωδικοποίησης βίντεο: H.264, H.265 και ProRes}

\subsubsection*{H.264 (AVC):}

\begin{itemize}
  \item \textbf{Χαρακτηριστικά}. Πιο διαδεδο$\mu$ένο, εξαιρετική συ$\mu$βατότητα.
  \item \textbf{Πλεονεκτή$\mu$ατα}: Σχεδόν όλες οι συσκευές, $\mu$έτρια συ$\mu$πίεση.
  \item \textbf{Μειονεκτή$\mu$ατα}. Χα$\mu$ηλότερη συ$\mu$πίεση από νέα, $\mu$εγαλύτερο.
  \item \textbf{Σενάρια χρήσης}. Web sharing, καθη$\mu$ερινά.
\end{itemize}

\subsubsection*{H.265 (HEVC):}

\begin{itemize}
  \item \textbf{Χαρακτηριστικά}: αναβαθ$\mu$ισ$\mu$ένο H.264, \textasciitilde{}50\% καλύτερη συ$\mu$πίεση.
  \item \textbf{Πλεονεκτή$\mu$ατα}. Ίδια ποιότητα, $\mu$ικρότερο, υψηλότερες ανάλυσεις και καρέ.
  \item \textbf{Μειονεκτή$\mu$ατα}. Encoding/decoding θέλει CPU, παλιά ασύ$\mu$βατα.
  \item \textbf{Σενάρια χρήσης}: 4K γυρίσ$\mu$ατα, περιορισ$\mu$ένος χώρος.
\end{itemize}

\subsubsection*{Apple ProRes:}

\begin{itemize}
  \item \textbf{Χαρακτηριστικά}. Apple pro codec, σχεδόν lossless.
  \item \textbf{Πλεονεκτή$\mu$ατα}. Ποιότητα υψηλή, post περιθώριο, ακριβές χρώ$\mu$α.
  \item \textbf{Μειονεκτή$\mu$ατα}. Τεράστιο, πολλός χώρος, pro εξοπλισ$\mu$ός.
  \item \textbf{Σενάρια χρήσης}: επαγγελ$\mu$ατικό βίντεο, κινη$\mu$ατογραφικό.
\end{itemize}

\subsubsection*{ProRes RAW:}

\begin{itemize}
  \item \textbf{Χαρακτηριστικά}. Συνδυάζει ProRes και RAW.
  \item \textbf{Τι είναι το RAW}. Ακατέργαστα sensor data, $\mu$έγιστη πληροφορία.
  \item \textbf{Πλεονεκτή$\mu$ατα}: Μεγάλο περιθώριο post, αλλάζει έκθεση, white balance χωρίς απώλεια.
  \item \textbf{Μειονεκτή$\mu$ατα}. Τεράστιο αρχείο, pro software.
  \item \textbf{Σενάρια χρήσης}:Επαγγελ$\mu$ατική παραγωγή υψηλού επιπέδου. Από τη σειρά iPhone 15 Pro και $\mu$ετά υποστηρίζεται εγγραφή βίντεο Apple ProRes (έως 4K 60fps), το iPhone 17 Pro υποστηρίζει περαιτέρω ProRes RAW απευθείας στην κά$\mu$ερα.
\end{itemize}

\booksubsection{III. HDR: τεχνολογία υψηλού δυνα$\mu$ικού εύρους}

Η τεχνολογία HDR (High Dynamic Range) διατηρεί ταυτόχρονα λεπτο$\mu$έρειες στις πιο φωτεινές και σκοτεινές περιοχές.

\begin{itemize}
  \item \textbf{HDR10}. Βασικό HDR, καλή συ$\mu$βατότητα.
  \item \textbf{HDR10+}: Samsung, dynamic metadata.
  \item \textbf{Dolby Vision}. Dolby top HDR, καλύτερη ποιότητα, $\mu$ε license.
  \item \textbf{HLG}: HDR broadcast, συ$\mu$βατό $\mu$ε κανονικές οθόνες.
\end{itemize}

Στα βίντεο HDR, χρειάζεστε οθόνη HDR για την πλήρη εικόνα.

\booksubsection{IV. Χρω$\mu$ατικός χώρος και βάθος χρώ$\mu$ατος}

\subsubsection*{Χρω$\mu$ατικός χώρος:}

\begin{itemize}
  \item \textbf{sRGB}. Standard, καλύτερη συ$\mu$βατότητα.
  \item \textbf{Display P3}. Apple wide gamut, πλουσιότερα χρώ$\mu$ατα.
  \item \textbf{Rec.2020}: χρω$\mu$ατικός χώρος HDR βίντεο, ευρύτερο gamut.
\end{itemize}

\subsubsection*{Βάθος χρώ$\mu$ατος:}

\begin{itemize}
  \item \textbf{8-bit}: 16,7 εκατ. χρώ$\mu$ατα, τυπικό βίντεο.
  \item \textbf{10-bit}: 1,07 δισ. χρώ$\mu$ατα, επαγγελ$\mu$ατικό, ο$\mu$αλές $\mu$εταβάσεις.
  \item \textbf{12-bit}: 68 δισ. χρώ$\mu$ατα, υψηλής ποιότητας.
\end{itemize}

\begin{tcbinfo}{Πρακτικές συ$\mu$βουλές}
\begin{itemize}
  \item Καθη$\mu$ερινά γυρίσ$\mu$ατα: 1080P/4K + 30fps + H.265, ισορροπία ποιότητας και $\mu$εγέθους.
  \item Αθλητικά γυρίσ$\mu$ατα: 1080P/4K + 60fps, για ο$\mu$αλότητα.
  \item Επαγγελ$\mu$ατική δη$\mu$ιουργία: 4K + ProRes (αν υποστηρίζεται), $\mu$έγιστο περιθώριο post.
  \item Συ$\mu$βουλή αποθήκευσης: τα βίντεο υψηλού bitrate καταλα$\mu$βάνουν πολύ χώρο, συνιστάται τακτική δη$\mu$ιουργία αντιγράφων ασφαλείας σε υπολογιστή ή cloud.
\end{itemize}

\end{tcbinfo}

\begin{tcbpractice}{Εξερευνήστε: Οι ρυθ$\mu$ίσεις βίντεο του κινητού σας}
\begin{enumerate}
  \item Ανοίξτε την εφαρ$\mu$ογή Κά$\mu$ερα, $\mu$εταβείτε στη λειτουργία βίντεο.
  \item Δείτε επιλογές ανάλυσης και καρέ, δοκι$\mu$άστε συνδυασ$\mu$ούς.
  \item Εάν υποστηρίζεται, ελέγξτε για προηγ$\mu$ένες επιλογές όπως HDR, ProRes.
  \item Γυρίστε βίντεο $\mu$ε διαφορετικές ρυθ$\mu$ίσεις, συγκρίνετε ποιότητα και $\mu$έγεθος.
\end{enumerate}

\end{tcbpractice}

\booksubsection{Πέ$\mu$πτο: Σύγκριση RAW $\mu$ε HEIF/JPEG}

\subsubsection*{Τι είναι η $\mu$ορφή RAW:}

Η $\mu$ορφή RAW είναι τα «ακατέργαστα δεδο$\mu$ένα» που καταγράφει ο sensor της κά$\mu$ερας — $\mu$η συ$\mu$πιεσ$\mu$ένες και ακατέργαστες πληροφορίες. Είναι σαν τα φρέσκα υλικά $\mu$αγειρικής, χρειάζονται post-processing για να γίνουν έτοι$\mu$α. Τα κινητά λόγω περιορισ$\mu$ών αποθήκευσης και επεξεργαστικής ισχύος τραβούν από προεπιλογή «έτοι$\mu$α» JPEG ή HEIF, το RAW χρειάζεται ειδική ενεργοποίηση.

\subsubsection*{Σύγκριση RAW vs HEIF vs JPEG:}

\begin{itemize}
  \item \textbf{JPEG}:Η παλαιότερη $\mu$ορφή, $\mu$ε την καλύτερη συ$\mu$βατότητα (ανοίγει σε οποιαδήποτε συσκευή). 8-bit color depth, lossy compression, $\mu$ικρότερο $\mu$έγεθος αρχείου (2-5MB ανά φωτογραφία). Περιορισ$\mu$ένος χώρος post-processing, η βίαιη φωτεινότητα σκοτεινών ση$\mu$είων δη$\mu$ιουργεί color banding.
  \item \textbf{HEIF/HEIC}:Νέα $\mu$ορφή που προωθεί η Apple, $\mu$ε το ίδιο επίπεδο ποιότητας το αρχείο είναι το $\mu$ισό του JPEG (1-3MB ανά φωτογραφία). Υποστηρίζει 10-bit color depth και alpha channel. Μειονέκτη$\mu$α: ορισ$\mu$ένες παλιές συσκευές και Windows χρειάζονται plugin για να ανοίξουν.
  \item \textbf{RAW}:Ασυ$\mu$πίεστα πρωτογενή δεδο$\mu$ένα, $\mu$εγαλύτερο $\mu$έγεθος αρχείου (20-100MB ανά φωτογραφία). Color depth συνήθως 12-16bit, $\mu$έγιστος χώρος post-processing. Μειονέκτη$\mu$α: απαιτεί επαγγελ$\mu$ατικό λογισ$\mu$ικό (Lightroom, Snapseed κ.ά.) για άνοιγ$\mu$α και επεξεργασία, δεν $\mu$πορεί να κοινοποιηθεί απευθείας.
\end{itemize}

\subsubsection*{Γιατί οι επαγγελ$\mu$ατίες φωτογράφοι βγάζουν RAW:}

\begin{itemize}
  \item \textbf{Η ισορροπία λευκού προσαρ$\mu$όζεται απωλεστικά στη $\mu$ετα-επεξεργασία}:Αν το JPEG έχει λάθος λευκή ισορροπία (π.χ. κιτρινωπό κάτω από λά$\mu$πα πυρακτώσεως), η διόρθωση post-processing $\mu$ειώνει την ποιότητα· το RAW $\mu$πορεί να επαναφέρει τη θερ$\mu$οκρασία χρώ$\mu$ατος εντελώς.
  \item \textbf{Μεγάλη ανοχή έκθεσης}: Στο JPEG, οι σκοτεινές περιοχές $\mu$ε υποέκθεση 2 σταθ$\mu$ών είναι σχεδόν $\mu$αύρες, και $\mu$ετά τη φωτίσηνονται γε$\mu$άτες θόρυβο. Το RAW $\mu$πορεί να "διασώσει" χρήσι$\mu$ες λεπτο$\mu$έρειες από σκοτεινές περιοχές $\mu$ε υποέκθεση 3-4 σταθ$\mu$ών. Το ίδιο ισχύει για την υπερέκθεση.
  \item \textbf{Πιο πλούσιες λεπτο$\mu$έρειες σκιών και φωτεινών ση$\mu$είων}: Το δυνα$\mu$ικό εύρος του RAW είναι πολύ $\mu$εγαλύτερο από του JPEG, $\mu$ε πολύ περισσότερο περιθώριο ώστε οι φωτεινές περιοχές να $\mu$ην υπερεκτίθενται και οι σκοτεινές να $\mu$ην είναι εντελώς $\mu$αύρες.
  \item \textbf{Όξυνση, $\mu$είωση θορύβου και διόρθωση φακού ελέγχονται στη $\mu$ετα-επεξεργασία}:Η επεξεργασία στην κά$\mu$ερα είναι «$\mu$ονής χρήσης», το RAW $\mu$πορεί να ρυθ$\mu$ιστεί επανειλη$\mu$$\mu$ένα, κάθε φορά από τα πρωτογενή δεδο$\mu$ένα.
\end{itemize}

\subsubsection*{Πώς να ενεργοποιήσετε RAW στο κινητό:}

\begin{itemize}
  \item \textbf{iPhone ProRAW}:Υποστηρίζεται $\mu$όνο σε iPhone 12 Pro και άνω Pro $\mu$οντέλα. Στις «Ρυθ$\mu$ίσεις Κά$\mu$ερα Μορφή» ενεργοποιήστε το «Apple ProRAW». Στη διεπαφή της κά$\mu$ερας θα ε$\mu$φανιστεί ένα κου$\mu$πί «RAW», πατήστε για ενεργοποίηση. Μέγεθος αρχείου περίπου 25MB/φωτογραφία, Συστήνεται να το ανοίγετε όταν χρειάζεστε.
  \item \textbf{Android Motion Raw / Ultra RAW}:Κάθε brand το ονο$\mu$άζει διαφορετικά. Σε Xiaomi/OPPO/vivo flagship, στον «Επαγγελ$\mu$ατικό λειτουργία» ή «Περισσότερα» της κά$\mu$ερας θα βρείτε τον διακόπτη RAW· στη Huawei στον «Επαγγελ$\mu$ατικό λειτουργία»· στη Samsung Galaxy S στον «Επαγγελ$\mu$ατικό» λειτουργία. Μέγεθος αρχείου συνήθως 20-50MB.
  \item \textbf{Γενικές συ$\mu$βουλές}:Στην καθη$\mu$ερινή λήψη δεν χρειάζεται RAW, το αρχείο είναι πολύ $\mu$εγάλο· ανοίξτε το $\mu$όνο σε ση$\mu$αντικές περιπτώσεις (τοπίο, πορτρέτο, post-processing).
\end{itemize}

\booksubsection{Έκτο: Αρχές υπολογιστικής φωτογραφίας}

Η σύγχρονη φωτογραφία κινητού δεν $\mu$πορεί πια να από την «υπολογιστική φωτογραφία» — δηλαδή τη βασίζεται σε αλγορίθ$\mu$ους και όχι $\mu$όνο σε φακούς και sensors για να παράγει φωτογραφίες. Η κατανόηση αυτών των βασικών αρχών σας βοηθά να καταλάβετε τι κάνει το τηλέφωνο τη στιγ$\mu$ή της λήψης.

\subsubsection*{Σύνθεση πολλαπλών καρέ:}

Όταν πατάτε το κλείστρο $\mu$ια φορά, το τηλέφωνο πραγ$\mu$ατικά τραβάει πολλές φωτογραφίες. Ο αλγόριθ$\mu$ος ευθυγρα$\mu$$\mu$ίζει και συγχωνεύει πολλαπλά καρέ για να παράγει $\mu$ια υψηλής ποιότητας φωτογραφία.

\begin{itemize}
  \item \textbf{Λειτουργία νυχτερινής λήψης}:Λα$\mu$βάνει πολλαπλά underexposed καρέ $\mu$έσα σε 1-3 δευτερόλεπτα, τα ευθυγρα$\mu$$\mu$ίζει και τα στοιβάζει — ο θόρυβος αλληλεταιρείται και η φωτεινότητα συσσωρεύεται. Αποτέλεσ$\mu$α: καθαρή εικόνα ακό$\mu$α και σε χα$\mu$ηλό φωτισ$\mu$ό.
  \item \textbf{HDR (Υψηλό Δυνα$\mu$ικό Εύρος)}:Λα$\mu$βάνει ξεχωριστά underexposed, κανονικές και overexposed φωτογραφίες (ή περισσότερες), συνθέτει $\mu$ια $\mu$ε λεπτο$\mu$έρειες σε φωτεινά και σκοτεινά ση$\mu$εία. Ιδιαίτερα αποτελεσ$\mu$ατικό σε αντίθετο φωτισ$\mu$ό.
  \item \textbf{Deep Fusion (Apple) / Photonic Engine}:Για σκηνές χα$\mu$ηλού-$\mu$εσαίου φωτισ$\mu$ού, αναλύει ανά pixel και συγχωνεύει πολλαπλά καρέ, εστιάζοντας στη βελτίωση υφής και λεπτο$\mu$ερειών — πλεκτά υφάσ$\mu$ατα, τρίχες θα είναι πιο καθαρά.
\end{itemize}

\subsubsection*{Ση$\mu$ασιολογική κατάτ$\mu$ηση:}

Η AI δεν επεξεργάζεται όλη την εικόνα $\mu$ε τον ίδιο τρόπο, αλλά πρώτα «κατανοεί» τι υπάρχει — ουρανός, πρόσωπα, γρασίδι, κτήρια, ρούχα — και $\mu$ετά εφαρ$\mu$όζει διαφορετική στρατηγική σε κάθε περιοχή.

\begin{itemize}
  \item \textbf{Ουρανός}: Ενισχύεται ξεχωριστά το $\mu$πλε και $\mu$ειώνεται η φωτεινότητα για αποφυγή υπερέκθεσης.
  \item \textbf{Πρόσωπο}: Αναγνωρίζεται η περιοχή χρώ$\mu$ατος δέρ$\mu$ατος, φωτίζεται και εξο$\mu$αλύνεται ξεχωριστά, χωρίς να πειράζονται τα χαρακτηριστικά και οι λεπτο$\mu$έρειες $\mu$αλλιών.
  \item \textbf{Γκαζόν/Φύλλα}: Ενισχύεται ο κορεσ$\mu$ός του πράσινου και βελτιώνεται η όξυνση λεπτο$\mu$ερειών.
  \item \textbf{Θόλωση φόντου}: Στη λειτουργία πορτρέτου, η AI αναγνωρίζει το περίγρα$\mu$$\mu$α του θέ$\mu$ατος, θολώνει το φόντο και κρατάει το θέ$\mu$α ευκρινές.
\end{itemize}

\subsubsection*{Εφαρ$\mu$ογές βαθιάς $\mu$άθησης:}

\begin{itemize}
  \item \textbf{AI αναγνώριση σκηνής}:Η κά$\mu$ερα αναγνωρίζει αυτό$\mu$ατα σκηνές όπως «φαγητό», «κατοικίδιο», «νυχτερινή», «έγγραφο» και προσαρ$\mu$όζει τις παρα$\mu$έτρους. Π.χ. η λειτουργία φαγητού ενισχύει τους θερ$\mu$ούς τόνους και τον κορεσ$\mu$ό.
  \item \textbf{AI θόλωση πορτρέτου}:Οι φακοί του κινητού δεν $\mu$πορούν φυσικά να πετύχουν το ρηχό depth of field $\mu$ιας κά$\mu$ερας DSLR, αλλά $\mu$ε $\mu$οντέλα deep learning προβλέπει τα όρια του θέ$\mu$ατος και θολώνει το φόντο. Οι νεότεροι αλγόριθ$\mu$οι χειρίζονται σωστά τις τρίχες και τα άκρα γυαλιού.
  \item \textbf{AI αφαίρεση / AI επέκταση εικόνας}:Επιλέξτε περαστικούς ή αντικεί$\mu$ενα στην εικόνα, η AI «συ$\mu$πληρώνει» το καλυ$\mu$$\mu$ένο τ$\mu$ή$\mu$α βάσει του περιβάλλοντος. Αυτή είναι εφαρ$\mu$ογή της generative AI στη φωτογραφία.
\end{itemize}

\subsubsection*{Υπολογιστική φωτογραφία vs οπτική φωτογραφία:}

\begin{itemize}
  \item \textbf{Πλεονεκτή$\mu$ατα υπολογιστικής φωτογραφίας}:Με $\mu$ικρό $\mu$έγεθος επιτυγχάνεται καλή ποιότητα, σε ακραίες σκηνές χα$\mu$ηλού/αντίθετου φωτισ$\mu$ού βγαίνουν καλές λήψεις, οι απλοί χρήστες γνώσεις φωτογραφίας για να βγάλουν «ό$\mu$ορφες» φωτογραφίες.
  \item \textbf{Μειονεκτή$\mu$ατα υπολογιστικής φωτογραφίας}:Η εικόνα $\mu$ερικές φορές «υπερ επεξεργάζεται» — υπερβολικό sharpening, χρω$\mu$ατική παρα$\mu$όρφωση, αφύσικος ουρανός· η σε κινού$\mu$ενο θέ$\mu$α αφήνει ghosting· οι λεπτο$\mu$έρειες είναι «υπολογισ$\mu$ένες» όχι πραγ$\mu$ατικά καταγραφό$\mu$ενες.
  \item \textbf{Πλεονεκτή$\mu$ατα οπτικής φωτογραφίας}:Μεγάλος sensor + $\mu$εγάλο διάφραγ$\mu$α φακού πραγ$\mu$ατικά το φως, φυσική ποιότητα, $\mu$εγάλος χώρος post-processing· δεν παράγει artifacts λόγω «$\mu$αντεψιάς» του αλγορίθ$\mu$ου.
  \item \textbf{Τάση}:Τα σύγχρονα flagship κινητά συνδυάζουν «οπτική + υπολογιστική» — $\mu$εγάλος sensor για πραγ$\mu$ατικά δεδο$\mu$ένα, αλγόριθ$\mu$ος για βελτιστοποιη$\mu$ένη παρουσίαση. Δεν είναι αντίθετα.
\end{itemize}

\booksubsection{Έβδο$\mu$ο: Τύποι φακών και εστιακά $\mu$ήκη}

Τα σύγχρονα flagship κινητά συνήθως έχουν τρεις έως τέσσερις φακούς, κάθε ένας για διαφορετική εστιακή απόσταση. Η κατανόηση των εστιακών αποστάσεων σας βοηθά να επιλέξετε τον σωστό φακό.

\subsubsection*{Συνηθισ$\mu$ένοι φακοί και ισοδύνα$\mu$η εστιακή απόσταση:}

\begin{itemize}
  \item \textbf{Υπερ$\mu$εγάλη γωνία (ισοδύνα$\mu$ο 13-24mm)}:Το πιο ευρύ οπτικό πεδίο, κατάλληλο για τοπίο, κτήρια, εσωτερικούς χώρους. Αποτυπώνει $\mu$εγάλες σκηνές, αλλά οι άκρες έχουν παρα$\mu$όρφωση (οι ευθείες γρα$\mu$$\mu$ές κα$\mu$πυλώνονται), στα πορτρέτα το πρόσωπο παρα$\mu$ορφώνεται. Κατάλληλο για «$\mu$εγάλες σκηνές», όχι για «close-up».
  \item \textbf{Κύριος φακός (ισοδύνα$\mu$ο 24-28mm)}:Ο πιο συχνά χρησι$\mu$οποιού$\mu$ενος φακός, $\mu$ε την καλύτερη ποιότητα και το $\mu$εγαλύτερο διάφραγ$\mu$α. Ισορροπεί οπτικό πεδίο και ποιότητα, για πορτρέτα και τοπία. Το «1x» των περισσότερων κινητών είναι η κύρια κά$\mu$ερα.
  \item \textbf{Τηλεφακός (ισοδύνα$\mu$ο 70-200mm)}:Για $\mu$ακρινά αντικεί$\mu$ενα και πορτρέτα close-up. Ο τηλεφακός συ$\mu$πιέζει τον χώρο (τα $\mu$ακρινά βουνά και το κοντινό άτο$\mu$ο φαίνονται πιο κοντά), στα πορτρέτα δίνει πιο φυσικές αναλογίες προσώπου και θόλωση φόντου.
  \item \textbf{Τηλεφακός periscope (ισοδύνα$\mu$ο 200mm+)}:Με ανακλαστικό σχεδιασ$\mu$ό, πετυχαίνει $\mu$εγαλύτερη εστιακή απόσταση $\mu$έσα στο πάχος του κινητού. Μπορεί να τραβήξει καθαρές λεπτο$\mu$έρειες από $\mu$εγάλη απόσταση, κατάλληλο για συναυλίες, αθλητικά, άγρια ζώα. Μόνο σε flagship $\mu$οντέλα.
\end{itemize}

\subsubsection*{Ανταλλαγή $\mu$εταξύ $\mu$εγέθους αισθητήρα και $\mu$εγέθους pixel:}

Ο sensor είναι η «βάση», καθορίζει πόσο φως $\mu$πορεί να λάβει. Στην οικονο$\mu$ία της φωτογραφίας «$\mu$εγαλύτερο sensor = καλύτερη ποιότητα» — $\mu$ε τον ίδιο αριθ$\mu$ό pixels, όσο $\mu$εγαλύτερος ο sensor, τόσο $\mu$εγαλύτερο κάθε pixel, τόσο περισσότερο φως, τόσο καλύτερη ποιότητα.

\begin{itemize}
  \item \textbf{Αισθητήρας $\mu$εγάλου φορ$\mu$ά}:Μεγάλη ποσότητα φωτός, καλή ποιότητα νυχτερινής λήψης, ευρύ dynamic range. Μειονέκτη$\mu$α: το τηλέφωνο πρέπει να είναι παχύτερο (ο φακός προεξέχει), υψηλότερο κόστος.
  \item \textbf{Υψηλά pixel $\mu$ικρού φορ$\mu$ά}:Όταν υπάρχει επαρκές φως τη $\mu$έρα, τα υψηλά pixels βγάζουν πιο καθαρές φωτογραφίες, διευκολύνοντας post-processing περικοπής. Αλλά σε χα$\mu$ηλό φως κάθε pixel δέχεται ανεπαρκές φως, $\mu$ε περισσότερο θόρυβο.
  \item \textbf{Επιλογές κύριων ναυαρχίδων}:Η κύρια κά$\mu$ερα είναι συνήθως 1/1.3 ίντσες έως 1 ίντσα (ορισ$\mu$ένα $\mu$οντέλα χρησι$\mu$οποιούν ήδη 1 ίντσα), $\mu$ε περίπου 50 megapixels — αυτό είναι το ση$\mu$είο ισορροπίας ποιότητας και $\mu$εγέθους.
\end{itemize}

\subsubsection*{Τεχνολογία συνένωσης pixel (Pixel Binning):}

Για να συνδυάσει τα υψηλά pixels τη $\mu$έρα και τα $\mu$εγάλα pixels τη νύχτα, οι κατασκευαστές χρησι$\mu$οποιούν «pixel binning» — συγχωνεύουν πολλαπλά $\mu$ικρά pixels σε ένα $\mu$εγάλο.

\begin{itemize}
  \item \textbf{Τετραπλή συνένωση pixel (Quad Bayer)}:Sensor 50 megapixels, εξάγει φωτογραφία 12,5 megapixels. Κάθε pixel εξόδου συντίθεται από 4 γειτονικά pixels, η ποσότητα φωτός αυξάνεται 4 φορές, πιο καθαρή ποιότητα.
  \item \textbf{Εννιαπλή συνένωση pixel (Nonacell)}:Για sensors άνω των 100 megapixels, εξάγει περίπου 11-12 megapixels. Ίδια αρχή, $\mu$εγαλύτερη ποσότητα φωτός.
  \item \textbf{Πραγ$\mu$ατικό αποτέλεσ$\mu$α}:Στην προεπιλεγ$\mu$ένη λειτουργία λήψης, αυτό που παίρνετε είναι η συγχωνευ$\mu$ένη φωτογραφία χα$\mu$ηλών pixels — καλή ποιότητα, $\mu$ικρό αρχείο. Για ακραίες λεπτο$\mu$έρειες (έγγραφα, τηλεφωτογραφία), $\mu$πορείτε να αλλάξετε χειροκίνητα σε «λειτουργία υψηλών pixels» για πλήρη ανάλυση.
\end{itemize}

\booksubsection{Όγδοο: Ειδικές λειτουργίες λήψης}

Εκτός από την καθη$\mu$ερινή λήψη, υπάρχουν και άλλες ειδικές λειτουργίες που, αν χρησι$\mu$οποιηθούν σωστά, βγάζουν εικόνες που δεν $\mu$πορείτε να πάρετε αλλιώς.

\subsubsection*{Λειτουργία πανορά$\mu$ατος:}

\begin{itemize}
  \item Κατά τη λήψη, ακολουθήστε το βελάκι και $\mu$ετακινήστε αργά και ο$\mu$αλά οριζόντια το τηλέφωνο, $\mu$ην βιαστείτε — αλλιώς θα ε$\mu$φανιστούν σφάλ$\mu$ατα ευθυγρά$\mu$$\mu$ισης.
  \item Κρατήστε το κινητό σταθερό σε κατακόρυφη κατεύθυνση και $\mu$ετακινήστε το $\mu$όνο οριζόντια. Αν κουνηθεί πάνω-κάτω, ο ορίζοντας θα γίνει κυ$\mu$ατιστός.
  \item Προσοχή όταν υπάρχουν κινού$\mu$ενα άτο$\mu$α ή οχή$\mu$ατα — κατά τη ραφή $\mu$πορεί να ε$\mu$φανιστούν «$\mu$ισά σώ$\mu$ατα» ή «φαντάσ$\mu$ατα».
  \item Το πανόρα$\mu$α $\mu$πορεί να βγει και κάθετα — κάθετη λήψη $\mu$ε οριζόντια ραφή, ιδανικό για ψηλά κτίρια ή καταρράκτες.
\end{itemize}

\subsubsection*{Καταβράδυνση και time-lapse:}

\begin{itemize}
  \item \textbf{Καταβράδυνση}:Λήψη $\mu$ε υψηλό frame rate 120fps ή 240fps, στην αναπαραγωγή γίνεται slow motion. Κατάλληλο για σταγόνες νερού, αθλητικές στιγ$\mu$ές, άλ$\mu$ατα κατοικίδιων. Απαιτείται επαρκές φως, η λειτουργία slow motion έχει $\mu$ικρό χρόνο έκθεσης, σε χα$\mu$ηλό φως η ποιότητα είναι κακή.
  \item \textbf{Time-lapse}:Λήψη πολλαπλών φωτογραφιών σε διαστή$\mu$ατα, συνθέτει βίντεο fast-forward. Κατάλληλο για ροή αυτοκινήτων, κίνηση σύννεφων, ηλιοβασίλε$\mu$α, άνθιση λουλουδιών. Διάρκεια λήψης από λίγα λεπτά έως ώρες, χρησι$\mu$οποιήστε υποχρεωτικά τρίποδο.
  \item \textbf{Μικρή τεχνική}:Στη λήψη time-lapse σύννεφων, στηρίξτε το τηλέφωνο, τραβήξτε 15-30 λεπτά, στο τέλος παίρνετε βίντεο περίπου 10 δευτερολέπτων, εντυπωσιακό αποτέλεσ$\mu$α.
\end{itemize}

\subsubsection*{Λειτουργία νυχτερινού ουρανού/αστρονο$\mu$ίας:}

\begin{itemize}
  \item Πρέπει να απο$\mu$ακρυνθείτε από τη φωτορύπανση της πόλης και να βρείτε προάστια ή ορεινές περιοχές χωρίς καθόλου παρε$\mu$βολές από φώτα δρό$\mu$ου. Στην πόλη δεν $\mu$πορείτε να φωτογραφίσετε τον Γαλαξία.
  \item Απαραίτητο τρίποδο, $\mu$ε το χέρι θα βγει θολή. Το κινητό πρέπει να σταθεροποιηθεί για $\mu$ερικά δευτερόλεπτα έως δεκάδες δευτερόλεπτα.
  \item Ο καιρός να είναι καθαρός χωρίς σύννεφα, το φεγγάρι να $\mu$ην είναι πολύ $\mu$εγάλο (πανσέληνος κρύβει τα αστέρια). Τα σκοτεινά βράδια αρχή ή τέλος σεληνιακού $\mu$ήνα είναι πιο κατάλληλα.
  \item Το «Νυχτερινό λειτουργία» του iPhone αναγνωρίζει αυτό$\mu$ατα τον έναστρο ουρανό και παρατείνει τον χρόνο έκθεσης (έως 30 δευτερόλεπτα)· το «Star Mode» ή «Astronomy Mode» των Android flagship έχει παρό$\mu$οιο αποτέλεσ$\mu$α. Κάντε υπο$\mu$ονή, $\mu$ετά το κλείστρο $\mu$ην κινήσετε το τηλέφωνο.
\end{itemize}

\subsubsection*{Λειτουργία macro:}

\begin{itemize}
  \item Για λήψη λουλουδιών, εντό$\mu$ων, κοσ$\mu$η$\mu$άτων, κει$\mu$ένου και άλλων $\mu$ικρών αντικει$\mu$ένων. Μπορεί να εστιάσει από απόσταση 1-2 εκατοστών.
  \item Ορισ$\mu$ένα τηλέφωνα έχουν ειδικό φακό macro (συνήθως 2 ή 5 megapixel $\mu$ικρού sensor), $\mu$έτριας ποιότητας· τα flagship χρησι$\mu$οποιούν τον ultra-wide φακό που κάνει και macro, $\mu$ε καλύτερη ποιότητα.
  \item Σε κοντινή λήψη ο φακός είναι πολύ κοντά στο θέ$\mu$α και $\mu$πλοκάρει το φως, χρησι$\mu$οποιήστε αν χρειάζεται lamp ή φακό για συ$\mu$πληρω$\mu$ατικό φωτισ$\mu$ό.
  \item Η απόσταση εστίασης είναι πολύ $\mu$ικρή, το depth of field πολύ ρηχό, η $\mu$ετακίνηση του θέ$\mu$ατος κατά λίγα χιλιοστά $\mu$πορεί να αποσυντονίσει. Κρατήστε σταθερό χέρι, ιδανικά το θέ$\mu$α να $\mu$ην κινείται.
\end{itemize}

\begin{tcbinfo}{Πρακτικές συ$\mu$βουλές στατικής φωτογραφίας}
\begin{itemize}
  \item Καθη$\mu$ερινή λήψη: HEIF/JPEG από προεπιλογή αρκεί, $\mu$ικρό $\mu$έγεθος αρχείου, επαρκής ποιότητα.
  \item Σε ση$\mu$αντικές περιπτώσεις (τοπίο, πορτρέτο, post-processing): ενεργοποιήστε RAW, διατηρήστε $\mu$έγιστο χώρο post-processing.
  \item Επιλέξτε τον σωστό φακό: τοπίο $\mu$ε υπερ$\mu$εγάλη γωνία, πορτρέτο $\mu$ε τηλεφακό, καθη$\mu$ερινά $\mu$ε τον κύριο φακό.
  \item Νυχτερινές σκηνές και αστέρια: σταθεροποιήστε το κινητό, ενεργοποιήστε λειτουργία νυχτερινής λήψης ή αστρονο$\mu$ίας.
  \item Τα τηλέφωνα $\mu$ε $\mu$εγάλο sensor (π.χ. 1 ίντσας) έχουν υψηλότερο όριο ποιότητας, αλλά οι συνήθειες χρήσης παρα$\mu$ένουν κρίσι$\mu$ες.
\end{itemize}

\end{tcbinfo}

\begin{tcbpractice}{Δοκι$\mu$άστε το: Εξερευνήστε τις δυνατότητες της κά$\mu$ερας σας}
\begin{enumerate}
  \item Ελέγξτε αν το τηλέφωνό σας υποστηρίζει λήψη RAW (Ρυθ$\mu$ίσειςΚά$\mu$εραΜορφή), δοκι$\mu$άστε να το ενεργοποιήσετε και να τραβήξετε $\mu$ια φωτογραφία, συγκρίνοντας το $\mu$έγεθος των αρχείων JPEG και RAW στην ίδια σκηνή.
  \item Εναλλάξτε στους τρεις φακούς ultra-wide, κύρια κά$\mu$ερα και τηλεφακό, τραβήξτε το ίδιο $\mu$ακρινό θέ$\mu$α, για να νιώσετε τη διαφορά οπτικής γωνίας και ποιότητας.
  \item Βρείτε $\mu$ια σκηνή $\mu$ε νυχτερινό λειτουργία ή star mode, σταθεροποιήστε το τηλέφωνο και τραβήξτε $\mu$ια νυχτερινή φωτογραφία.
  \item Δοκι$\mu$άστε να τραβήξετε $\mu$ια $\mu$εγάλη σκηνή $\mu$ε πανόρα$\mu$α οριζόντια, προσέχοντας να $\mu$ετακινείτε το κινητό ο$\mu$αλά.
\end{enumerate}

\end{tcbpractice}
